% !TeX encoding = UTF-8
% !TeX spellcheck = en_GB
% !TeX root = ../thesis.tex


\chapter{Introduction}
\label{chapter:introduction}
From online banking to social communications and e-commerce, web applications have become the backbone of modern digital life. A central aspect to the security of these systems is the concept of a web session, which is a temporary authenticated connection between a user and a server established after a successful login. Such sessions have an expiry date and are managed through cookies, which the browser presents in every request to prove the user's identity. While this mechanism is widely adopted and convenient, it introduces a vulnerability: if the attacker can obtain a valid session token, they can impersonate a legitimate user without ever needing to use their credentials. This attack is known as Session Hijacking. It is one of the most prevalent threats to web security. 

Session hijacking commonly occurs through stealing session cookies. Although recent improvements to the browser handling of cookies have made cookie theft more complicated, the cookies can still be stolen via malware \cite{shen2021phishing}. Once a session cookie is captured, the attacker has full access to the victim's account as long as the session remains valid. The server has no reliable mechanism to distinguish between the legitimate user and an impostor presenting the session cookie since they both look identical from the server's perspective.

Such session hijacking attacks also bypass any Multi-Factor Authentication (MFA) imposed by the service. MFA only protects the initial login step, but once the session is established and the token is produced, the session itself is no longer protected by any additional mechanism. This gap between login-time security and session-time security is a significant blind spot in most current web security systems which motivates the need for continuous session validation.

A promising method to address this gap is through continuous authentication (CA), where the user is continuously authenticated throughout the session rather than relying solely on the initial login. One of the ways by which CA can be achieved is hardware binding, where a web session is associated with some device or hardware in their possession. If the server can continuously verify that the requests are coming from a specific physical device, a stolen session token becomes useless to an attacker using a different machine. Physical Unclonable Functions (PUF) can be a compelling primitive for this purpose. A PUF exploits the unique variations introduced by the manufacturing process on every device. When challenged with some input, an ideal PUF produces a response unique to that device; such a response is impossible to clone or replicate, making it ideal to be used as a primitive for key generation or other cryptographic operations.


%The PUSH protocol proposed by \cite{10.1007/978-3-032-07894-0_7} demonstrated that PUF-based hardware binding is a viable approach to continuous session validation. While the PUSH protocol represents an important step towards PUF-based hardware binding, the proposed implementation comes with two notable limitations that constrain its applicability. First, to challenge SRAM PUFs, a power cycle is required and therefore a cool-down period is in order before the next challenge. This introduces significant latency that is incompatible with the low-latency demands of CA. Second, the SRAM PUF is a weak PUF, meaning it has a very limited number of Challenge-Response Pairs (CRPs) which introduces the reuse of keys.

The PUSH protocol proposed by \cite{10.1007/978-3-032-07894-0_7} demonstrated that PUF-based hardware binding is a viable approach to continuous session validation. However, PUSH's design fully trusts the host machine to faithfully mediate the interaction between the browser and the dongle. This assumption introduces the following flaw: cookie theft already requires the attacker's malware to have remote code execution on the host, and that same capability is enough for the malware to also act as a proxy to the dongle. By silently forwarding the server-sent signature requests to the dongle and getting the signed cookies, an attacker can maintain a hijacked session without the user noticing. Another limitation of PUSH is tied to the usage of an SRAM PUF. The SRAM PUF is a weak PUF, meaning it has a limited number of challenge-response pairs (CRPs). This forces keys to be reused once all the CRP space is exhausted, which is subpar for security.
 
Moreover, the use of PUFs in web security remains a very under-researched area. The vast majority of PUF literature focuses on IoT-related scenarios where PUFs are well established. Their application to web security, where the threat model and performance constraints are different, has comparatively received little attention. To our knowledge, PUSH is the only work in this direction. While PUSH uses a weak PUF, no other work has explored the usage of reconfigurable PUFs for CA.

This thesis aims to address both the limitations of the PUSH protocol and the broader research gap by proposing the REPUSH protocol (Reconfigurable PUF against Session Hijacking). REPUSH improves on PUSH from two sides. First, REPUSH drops the assumption that the host can be trusted. Only the dongle and the server are considered trusted parties, while the host, the browser and the dongle manager are considered as compromised. This stronger threat model motivates many protocol changes: the local dongle manager simply acts as a translator between the browser and dongle, and all cryptographic operations happen on the dongle and server. The requests are split into non-sensitive and sensitive where the latter require explicit user approval on the dongle's own touch screen. The dongle authenticates the server using a certificate authority (CA), and a new reconfiguration phase renews the session at the end of every session rather than reusing the same key indefinitely. Second, to address the limited CRP space, REPUSH replaces the SRAM PUF with a Logically Reconfigurable PUF (LR-PUF) \cite{10.1007/978-3-642-23951-9_25} built on top of a CHOICE-PUF \cite{9556423} implemented on a Field Programmable Gate Array board (FPGA). The LR-PUF wraps the underlying weak CHOICE-PUF with hashing and internal state logic to enable reconfiguration. The reconfiguration mechanism removes the CRP space as the binding constraint on the number of keys that can be generated. The dongle is made up of an FPGA, an Arduino and a touch screen. The FPGA hosts the CHOICE-PUF and the LR-PUF logic around it, whilst an Arduino board handles communication with the host, the cryptographic operations, and drives the dongle's touch screen display. This dongle serves as a proof-of-concept demonstrating that a reconfigurable hardware-bound session validation protocol can be built with accessible and relatively low-cost components. 

\section{Research Questions}
The limitations of the PUSH protocol and the lack of research addressing PUF-based web security give rise to a set of research questions that this thesis aims to answer. The first two research questions concern the protocol design itself: one addresses the protocol changes needed to respect the stronger threat model, and the other addresses the integration of an LR-PUF into the protocol. The third question examines the security advantages brought by the design changes. The fourth question examines the practical implementation of such a protocol. These four questions frame both the theoretical and practical contributions of this thesis:
\begin{itemize}
    \item \textbf{Research question 1:} How must the PUSH protocol be redesigned to operate under a threat model that assumes host compromise?
    \item \textbf{Research question 2:} How can the PUSH protocol be modified to accommodate a logically reconfigurable PUF, i.e, how should reconfiguration be integrated into the session validation life-cycle?
    \item \textbf{Research question 3:} What security improvements does the REPUSH protocol provide over the original PUSH protocol?
    \item \textbf{Research question 4:} What are the performance and security overheads introduced by a PUF-enabled hardware dongle (FPGA and micro-controller) in a continuous session validation protocol?
\end{itemize}

\section{Thesis Outline}
The remainder of this thesis is structured as follows. Chapter \ref{chapter:background} provides the necessary background on Session Hijacking, the PUSH protocol, Physical Unclonable Functions (PUFs), Choice-PUF, Logically Reconfigurable PUF (LR-PUF) and Field Programmable Gate Arrays (FPGAs). Chapter \ref{chapter:related_work} surveys related works on methods against Session Hijacking, Reconfigurable PUFs and implementations of PUFs on FPGAs. Chapter \ref{chapter:design} presents the architecture of the REPUSH protocol and the hardware dongle. Chapter \ref{chapter:implementation} describes the implementation of each party of the protocol. Chapter \ref{chapter:evaluation} evaluates the protocol's security and performance, and compares it against the original PUSH protocol. Finally, Chapter \ref{chapter:conclusions} concludes the thesis and proposes directions for future work.




%\lipsum[3-56]
